\documentclass[aspectratio=169]{beamer}
% The values that it accepts are: 1610, 169, 149, 54, 43 and 32, which stand for the ratios 16:10, 16:9, 14:9, 5:4, 4:3 and 3:2, respectively.

\newcommand{\sigle}{ABC-1234}
\title[Titre principal qui apparait au dessus des diapositives]{} 

\input{varia/preamble}


%\subtitle[]{}

\author[O. Côté]{\\ 
O. Côté}

\institute%[Université Laval]
{
 École d'actuariat \\
 Université Laval, Québec, Canada \\
 \medskip
 %{\emph{olivier.cote.12@ulaval.ca}}
}%\today} % \today will show current date. 


%% ==========================================
%%  Informations de publication
%%  (titre et al. dans couverture-avant.tex)
%% ==========================================
\renewcommand{\year}{2022}
\renewcommand{\month}{05}
\renewcommand{\day}{13}


\setcounter{tocdepth}{2}
\newcommand{\acc}[1]{\textcolor{ulred}{\textbf{#1}}}


\usepackage[nostamp]{draftwatermark}
\begin{document}



\sloppy

%%%% ==========
%%%  page titre
%%%% =========

% %% PAge titre (edition Vincent)
% \include{0_pagetitre}

% % Page titre conventionnelle
% \begin{frame}[label=titre, plain]
%  \titlepage
%  \begin{center}%\includegraphics[height=1.2cm]{COPL}%
%   %\hspace{2cm}
%   \includegraphics[height=2cm]{img/ul_logo.png}\end{center}
% \end{frame}

%%%% ==========
%%%  table des matières
%%%% =========

% \begin{frame}[label=toc]{Table des matières}
%  \setlength{\leftskip}{5cm}%
%  \tableofcontents[subsectionstyle=hide]
% \end{frame}

%%%% ==========
%%%  Contenu
%%%% =========

\begin{frame}[label=intro]{Introduction slide}
	This document shows how to use the \emph{ulaval} Beamer template.
	It gives examples to see how the resulting document looks like.
\end{frame}

\section{Main Section}

\subsection{List styles}

\subsubsection{Itemize}

\begin{frame}[label=itemize]\frametitle{Itemize sample} 
	\begin{itemize}
		\item Item 1
		\item Item 2
		\begin{itemize}
			\item Sub item 1
			\item Sub item 2
			\begin{itemize}
				\item Sub sub sub item 1
				\item Sub sub sub item 2
			\end{itemize}
			\item Sub item 3
		\end{itemize}
		\item Item 3
	\end{itemize}
\end{frame}

\subsubsection{Enumerate}

\begin{frame}[label=enumerate]\frametitle{Enumerate sample} 
	\begin{enumerate}
		\item Item 1
		\item Item 2
		\begin{enumerate}
			\item Sub item 1
			\item Sub item 2
			\begin{enumerate}
				\item Sub sub sub item 1
				\item Sub sub sub item 2
			\end{enumerate}
			\item Sub item 3
		\end{enumerate}
		\item Item 3
	\end{enumerate}
\end{frame}

\subsubsection{Description}

\begin{frame}[label=description]\frametitle{Description sample}
    
\begin{description}
	\item[Term 1:] Definition 1
	\item[Term 2:] Definition 2
	\item[Term 3:] Definition 3
\end{description}

\end{frame}

\subsection{Boxes styles}

\begin{frame}[label=boxes]\frametitle{Boxes Styles}
    
\begin{block}{Block Title}
	Block content
\end{block}

\begin{alertblock}{Alert Block Title}
	Alert block content
\end{alertblock}

\begin{exampleblock}{Example Block Title}
	Example block content
\end{exampleblock}

\end{frame}

\subsection{Block environments}

\begin{frame}[label=environments]\frametitle{Environments Samples}
    
	\begin{definition}
	Definition content
	\end{definition}
  
	\begin{example}
	Example content
	\end{example}

	\begin{proof}
	Proof content
	\end{proof}  
    
	\begin{theorem}
	Theorem content
	\end{theorem}

\end{frame}


\subsection{Math}

\begin{frame}[label=math]\frametitle{Math}
    
\begin{equation}
	V_0 = k_0 \rho \sqrt{n_1^2 - n_2^2}
\end{equation}

\end{frame}


\subsection{Text environments}

\begin{frame}[label=text]\frametitle{Text Environments}
    
\begin{quotation}
  Quotation environment line 1\\
  Quotation environment line 2
\end{quotation}
\begin{quote}
  Quote environment line 1\\
  Quote environment line 2
\end{quote}
\begin{semiverbatim}
  Semiverbatim environment
\end{semiverbatim}
\begin{verse}
  Verse environment line 1\\
  Verse environment line 2
\end{verse}  

\end{frame}

\section{Conclusion}

\begin{frame}[label=conclu]{Conclusion slide}
	That's all folks!
\end{frame}

% End of slides
\end{document}

